% Einheit 2 von 2 – Linearkombination mit Nebenbedingung (bank/linearkombination-und-lineare-abhaengigkeit/e2.jsonl, zusammenbau v0.1)
%% TODO Zeitmarke Einheit 2: Marken-Zeile nicht lesbar
%% TODO Zweigzeile Teil 1: was hier gelernt wird (Satz in Schülersprache aus dem Einheitstitel) – Einheit 2
\einheitenkopf[e2]{Einheit 2 von 2 · Linearkombination mit Nebenbedingung}
\zweigzeile{Abitur LK}
\setcounter{aufgabe}{7}

%% TODO Titel: Ich-kann-Satz fehlt in der Bank (Platzhalter: Kettenname) – Nr. 8
%% TODO Anweisung: ein Satz über der ersten Teilaufgabe fehlt in der Bank – Nr. 8
\begin{aufgabe}{Linearkombination mit Nebenbedingung}
%% TODO \rechenplatz{n}: Schrittzahl des Grundfalls fehlt in der Bank (nur bei fester Schreibform, 2.3 b)
\begin{teile}
\teil $\vec{OX} = 0{,}3 \cdot \vec{OA} + 0{,}7 \cdot \vec{OB}$. Ergeben die Koeffizienten zusammen $1$, und liegen beide zwischen $0$ und $1$? \\ \kreuz{Summe 1, beide in [0; 1]} \\ \kreuz{Summe 1, nicht beide in [0; 1]} \\ \kreuz{Summe nicht 1}
\teil Forme $\vec{OX} = p \cdot \vec{OA} + q \cdot \vec{OB}$ mit $p + q = 1$ für $A(1 | 0 | 2)$ und $B(3 | 4 | 0)$ in eine Parameterform um. $\vec{OX} =$ \_\_
\teil Für $A(2 | 0 | 1)$ und $B(4 | 2 | 5)$ sei $\vec{OX} = p \cdot \vec{OA} + q \cdot \vec{OB}$ mit $p + q = 1$ und $p, q \in [0; 1]$. Welcher Punkt ergibt sich für $q = 0$, für $q = 1$ und für $q = 0{,}5$?
\teil Ein Dachstuhl hat die Form einer Pyramide $ABCDS$ mit $B(6 | 0 | 0)$ und der Spitze $S(3 | 3 | 4)$. Für einen Punkt $X$ gilt $\vec{OX} = p \cdot \vec{OB} + q \cdot \vec{OS}$ mit $p, q \in [0; 1]$ und $p + q = 1$. Weise nach, dass $X$ auf der Kante $BS$ liegt. \hfill (Abitur 2017 LK)
\end{teile}
\end{aufgabe}

%% TODO Titel: Ich-kann-Satz fehlt in der Bank (Platzhalter: Kettenname) – Nr. 9
%% TODO Anweisung: ein Satz über der ersten Teilaufgabe fehlt in der Bank – Nr. 9
\begin{aufgabe}{Linearkombination mit Nebenbedingung – Fehler finden, Begründen}
\begin{teile}
\teil Nina soll zeigen, dass $X$ mit $\vec{OX} = p \cdot \vec{OA} + q \cdot \vec{OB}$, $p + q = 1$ und $p, q \in [0; 1]$ auf der Strecke $AB$ liegt, wobei $A(1 | 1 | 0)$ und $B(3 | 5 | 2)$. Sie rechnet nur $p = q = 0{,}5$ und findet den Mittelpunkt von $AB$. Finde den Fehler und führe den Nachweis.
\teil Begründe, warum man mit $p = 1 - q$ aus $p \cdot \vec{OA} + q \cdot \vec{OB}$ eine Parameterform erhält.
\end{teile}
\end{aufgabe}
